\section{What makes an acoustic stimulus biologically meaningful?}
\label{sec:q1}

% Part I — Recognition and ecological meaning (Q1).
% See docs/OUTLINE.md — Q1.
% Draft 1. Style model: Magrath (agents/STYLE.md).

A claim that a sound is biologically meaningful is a claim that the receiver responds to information the sound carries, not merely to its energy, novelty, or aversiveness. Testing that claim requires a control that preserves the physical properties a receiver might be responding to while removing the information. The comparison most often used, natural calls against synthetic or artificial sounds, rarely achieves this. In food-restricted rats, natural 50-kHz calls evoked approach and transient dopamine release in the nucleus accumbens, whereas natural 22-kHz calls and noise matched in duration and amplitude evoked neither \citep{WilluhnEtAl2014}. The result shows that rats distinguish these calls from matched noise, but noise and calls still differ in fine spectrotemporal structure, so it cannot identify which property of the 50-kHz call carries the effect.

Studies that simplify calls systematically indicate that naturalness is neither necessary nor sufficient for a response. Sine-wave approximations of 50-kHz calls, lacking natural frequency and amplitude modulation, still induced approach, particularly in adult rats; juveniles responded differently in magnitude, so the natural and simplified stimuli were not equivalent \citep{WohrSchwarting2007}. Comparing natural 22-kHz and 50-kHz calls with computer-generated tones, \citet{OlszynskiEtAl2020} found that natural calls usually produced larger vocal, cardiac, and locomotor responses, but the direction of each effect was broadly the same for natural and artificial stimuli, and single housing enlarged several effects. At the neural level, frequency-matched tones reproduced most of the amygdala discrimination between the two call types \citep{ParsanaEtAl2012a} (\secref{sec:q3}). These studies indicate that much of what distinguishes rat call types for a receiver is carried by frequency band and gross temporal structure, and that modulation adds to, rather than creates, the response. ``Natural'' is therefore not a single causal variable but a bundle of acoustic dimensions, only some of which matter for a given receiver and endpoint.

The advantage of a natural call can also depend on its physical intensity. Rats oriented more to a recorded distress squeal than to a synthetic mimic matched in burst timing at 100~dB, but not at 80~dB \citep{EvansHammond1983a}. Significance and intensity thus interact rather than adding independently, and a comparison at a single sound level can either reveal or conceal a difference between natural and synthetic stimuli. Intensity and aversiveness are, conversely, not evidence of meaning:
% VERIFY: MongeauEtAl2003 (not full-text checked): 17--20 kHz FM sweep at 85 dB; flight versus freezing.
a standardized synthetic frequency sweep, not a natural call, elicited flight or freezing in mice \citep{MongeauEtAl2003}. A defensive response to a sound shows that the sound is aversive in that context, not that the receiver has extracted information about a threat.

Synthetic stimuli can also acquire meaning through experience. Mouse pup calls cluster near an inter-syllable interval of about 175~ms, and maternal experience broadened and shifted both retrieval behavior and auditory cortical tuning to synthetic calls in which this interval was altered \citep{SchiavoEtAl2020}. The system combines an intrinsic bias toward species-typical timing with plasticity that extends meaning to acoustic variants the receiver has never heard from a pup. The result concerns maternal recognition of pup calls and should not be generalized to a template for all call recognition.

Field studies of sciurids identify which acoustic dimensions free-living receivers use to assess risk.
% VERIFY: BlumsteinArmitage1997 (not full-text checked): calls track risk rather than predator type; receivers use call rate.
Alarm calls of yellow-bellied marmots (\textit{Marmota flaviventer}) do not map onto predator species; instead, call rate and some acoustic parameters vary with risk, and receivers use call rate as information about it \citep{BlumsteinArmitage1997}.
% VERIFY: BlumsteinRecapet2009 (not full-text checked): nonlinear components added to marmot calls reduced subsequent foraging.
% VERIFY: BlumsteinRecapet2009: comparison was against unmanipulated calls (and calls with added silence).
Adding noisy nonlinear components to marmot alarm calls reduced subsequent foraging more than unmanipulated calls did, so particular spectral features also raise responsiveness \citep{BlumsteinRecapet2009}.
% VERIFY: SloanEtAl2005 and SwanHare2008 (not full-text checked): call-type differences in vigilance; alerting effect of primary syllables.
In Richardson's ground squirrels, chirps, whistles, and chucks alter vigilance differently \citep{SloanEtAl2005}, and the primary syllables of repeated calls have especially strong alerting effects \citep{SwanHare2008}.
Gunnison's prairie dogs (\textit{Cynomys gunnisoni}) provide the clearest case of predator-specific information: playback of alarm calls originally given to humans, domestic dogs, or coyotes reproduced the distinct escape responses that each live predator elicits, whereas a control sound matched in rhythm, frequency range, and amplitude did not \citep{KiriazisSlobodchikoff2006}. The authors interpret the calls as referential; the playback results establish that the calls carry information sufficient for predator-specific escape, which does not require that they function as lexical labels. Urgency coding in marmots and predator-specific information in prairie dogs are two solutions to the same problem, and neither has been shown to be general across sciurids.

Information about the source also contributes to meaning, independently of call structure. Marmots responded more to calls from juveniles than to calls from adult females, but did not distinguish familiar from unfamiliar callers \citep{BlumsteinDaniel2004}.
% VERIFY: WarrenEtAl2025 (not full-text checked): partner versus stranger call discrimination in female prairie voles.
Caller identity is carried by 50-kHz calls in rats \citep{VielleEtAl2021} and by vocalizations in prairie voles (\textit{Microtus ochrogaster}) \citep{WarrenEtAl2025}; we return to how receivers use identity in \secref{sec:q5}.
% VERIFY: SchmidtEtAl2008 (not full-text checked): chipmunk foraging responses to titmouse mobbing calls versus hawk calls.
The relevance of a sound cannot be ranked by whether its source is a predator: eastern chipmunks (\textit{Tamias striatus}) changed their foraging more in response to the mobbing calls of titmice than to the calls of a hawk \citep{SchmidtEtAl2008}.
An indirect social cue that a predator is present and active can thus outweigh the predator's own vocalization.

These results suggest what a convincing test of acoustic meaning requires. The control should be matched to the call along the dimensions that the receiver could plausibly use (level, spectrum, temporal envelope, and modulation), and the comparison should be repeated at more than one sound level and in receivers whose experience is known, because both intensity and experience change the outcome. Laboratory studies have approached this standard for rat ultrasonic calls and mouse pup calls; field studies supply the strongest evidence that call structure carries information about risk, but they have rarely compared natural calls with acoustically matched synthetic reconstructions. Combining the two remains the most direct way to establish which acoustic dimensions carry meaning for free-living receivers.
\tabref{tab:summary-q1} consolidates these claims for later condensation or figure drafting.

% Intermediate summary matrix for Q1. May later be condensed, moved to an appendix, or redrawn as a figure.
\begin{table}[tbp]
\centering
\footnotesize
\singlespacing
\caption{Working summary for Q1: what carries acoustic meaning. Status reflects the strength of evidence reviewed in this section, not a formal meta-analysis.}
\label{tab:summary-q1}
\begin{tabular}{@{}p{1.55in}p{3.7in}p{1.0in}@{}}
\toprule
\textbf{Claim} & \textbf{Finding} & \textbf{Status} \\
\midrule
Natural $>$ synthetic & Often true vs noise or poorly matched controls; fails when frequency or timing is matched & Mixed \\
\addlinespace[2pt]
Naturalness is causal & Frequency band and gross timing often suffice; modulation adds to, rather than creates, responses & Strong against \\
\addlinespace[2pt]
Significance $\times$ intensity & Natural-call advantage can appear at high SPL and vanish at lower SPL & Thin (1 study) \\
\addlinespace[2pt]
Synthetic can acquire meaning & Maternal experience retunes responses to pup-call interval morphs & Strong (maternal) \\
\addlinespace[2pt]
Field call structure & Urgency (marmot rate, nonlinear noise); predator-specific escape (prairie dogs); call type and syllable position (ground squirrels) & Strong \\
\addlinespace[2pt]
Source ranking & Indirect social cues can outweigh predator vocalizations & Caution \\
\addlinespace[2pt]
Adequate control & Match level, spectrum, envelope, and modulation; test $>$1 SPL; know receiver experience & Gap in field \\
\bottomrule
\end{tabular}
\end{table}

